\section{Conclusions}

This project set out to explore a question with real clinical stakes: can a deep sequence model pick up on the physiological differences between bipolar and unipolar depression purely from the way a person moves throughout their day? We proposed a hybrid 1D-CNN-LSTM architecture that processes raw wrist actigraphy without any hand-crafted feature engineering, reasoning that the combination of local convolutional pattern detection and long-range LSTM memory would be better matched to the actual structure of this signal than classical approaches.

Our two-stage experimental design reflects the layered difficulty of the problem. Experiment~1, the healthy-versus-depressed baseline, serves as a sanity check and provides a meaningful comparison point against prior work on the Depresjon dataset. If the model struggles here, it would suggest a fundamental architectural or data problem worth resolving before moving on. Experiment~2 is where the real challenge lies. Separating bipolar from unipolar depression based on motor activity alone is a task that clinicians themselves find difficult from behavioral observation; asking a model to do it from wrist accelerometry is ambitious. We do not claim this will be easy or that the results will be clean. What we do claim is that multi-day temporal modeling is a more principled approach to this problem than computing variance over fixed windows, and the experimental results should tell us whether that intuition is correct.

Beyond the immediate classification results, this work points toward a broader possibility: that passive, continuous monitoring from consumer-grade wearables could one day provide clinicians with an objective data trail that supplements and perhaps eventually improves the diagnostic process for mood disorders. The path from a promising F1-Score on a 55-person dataset to a clinically validated tool is of course long and uncertain. But it has to start somewhere, and demonstrating that the temporal structure in actigraphy carries diagnostic signal is a reasonable first step.
